% proof.tex -- the wall rings are *2 for every n (whole-board game).
% Standalone wrapper. The part inside \begin{wallproof} ... \end{wallproof} is typeset as a section of the supplement;
% it only uses the macros \lab, \nim, \cyl, \ins (fall-backs below).
% Labels: \lab{P} = proved by hand, \lab{C2} = checked by two independently written programs.
\documentclass[10pt]{article}
\usepackage{amsmath,amssymb,amsthm}
\usepackage[margin=2cm]{geometry}
\providecommand{\lab}[1]{\textsf{\footnotesize[#1]}}
\providecommand{\nim}[1]{\mathord{\star#1}}
\providecommand{\cyl}[1]{C_{#1}}
\providecommand{\ins}{\operatorname{ins}}
\newtheorem{theorem}{Theorem}
\newtheorem{lemma}{Lemma}
\newenvironment{wallproof}{}{}
\begin{document}

\begin{wallproof}
\subsection*{The wall rings are $\nim{2}$ for every $n$ (whole-board game)}

\paragraph{Model.}
The board $\cyl{n}$ has ranks $1,\dots,8$ and $n$ files, numbered modulo $n$, so file $n-1$ is next to file $0$. White pawns move
towards rank 8, Black pawns towards rank 1. A pawn steps one square forward to an empty square, or two squares from its home rank (2 for
White, 7 for Black) if both squares are empty, and captures an enemy pawn one square diagonally forward, also across the seam
between files $n-1$ and $0$. After a double step an \emph{e.p.\ flag} is recorded if, and only if, an enemy pawn stands beside the landing
square. The flag names the pawn that stepped; the opponent of that pawn holds the right to capture it en passant, with a pawn standing beside it, onto the skipped square.
Every pawn move removes the old flag (and may record a new one); a heap move leaves the board and the flag unchanged.
A pawn is \emph{passed} if no enemy pawn stands strictly ahead of it on its own file or on a neighbouring file, and it is not the
flagged pawn (a pawn that can still be taken e.p.\ is not passed).
A position reached by a pawn move is \emph{terminal} if some side has a passed pawn. If only one side has one, that side wins. If both
have one, the race rule ``count'' decides: the distance of a side is the least value of $8-r-[r=2]$ (White) or $r-1-[r=7]$ (Black) over
its passed pawns on rank $r$ that have no pawn of either colour ahead on their own file (the empty minimum is
$\infty$); the smaller distance wins, and a tie goes to the side that did not make the last move.
In the \emph{whole-board game} $G+\nim h$ the players alternate; the player to move makes a pawn move on $G$ or lowers the heap.
A pawn move to a terminal position ends the game, won by the winner of that position whoever would move next. A player with no move
loses. There is no splitting of $G$ into components.

Let $X$ be the player to move in a non-terminal position $p$ and $Y$ the other. A pawn move $\mu$ of $X$ to $x$ is \emph{T} (\emph{L}) if
$x$ is terminal and won by $X$ (by $Y$); \emph{P} (poisoned) if $x$ is not terminal and $Y$ has a pawn move from $x$ to a terminal
position won by $Y$; \emph{K} otherwise. The classes are defined for pawn moves only and depend on who moves. The \emph{reduced game}
$R(g)$ of a non-terminal position $g$ has as nodes the positions reachable from $g$ by K-moves (of either side, in any order); at a node,
each side has its K-moves (to nodes) and possibly T-moves. In $R(g)+\nim h$ the player to move may make a T-move (and wins), a K-move,
or a heap move.

\paragraph{The rings.}
A file is written by its pawns ($W23B4$: white pawns on ranks 2 and 3, a black pawn on rank 4); files are listed from $0$ to $n-1$.
$O_n$ ($n=2k+3\ge5$): $W23B4,\,B5,\,W2B5,\,(W3,B5)^k$.\quad
$E_n$ ($n=2m+6\ge6$): $B5,\,W2B5,\,W35B6,\,W2,\,(B4,W2)^{m+1}$. Both have $n+3$ pawns and no flag.
Put $G_n=O_n$ for odd $n\ge5$ and $G_n=E_n$ for even $n\ge6$, and $\mathrm{start}_n$ for this position.
A \emph{wall} is a run of consecutive files that start with single pawns, alternately a white pawn and a black pawn two ranks above it:
$W3/B5$ (type O) or $W2/B4$ (type E). In $G_n$ the wall is files $3,\dots,n-1$ (type O: $W3$ on odd, $B5$ on even files; type E: $W2$ on
odd, $B4$ on even files). A \emph{pattern file} is a wall file whose content is exactly its starting pawn (same colour, same rank) and
nothing else; its pawn is a \emph{wall pawn}. On a pattern file, each neighbouring pattern file holds a single enemy pawn two ranks
further on in the direction of the file's own pawn.

\begin{theorem}
For every $n\ge5$ and every heap $0\le h<256$, the whole-board position $G_n+\nim h$ is a P-position if $h=2$ and an N-position
otherwise. So $G_n$ and $\nim2$ have the same outcome in every sum with a nim-heap; in this sense $G_n=\nim2$.
\end{theorem}

\begin{lemma}[Reduction; \lab{P}]
Let $g$ be a non-terminal position, $h\ge0$, and let either side move first. That side wins $g+\nim h$ iff it wins $R(g)+\nim h$.
\end{lemma}
\begin{proof}
Let $\pi(p)$ be the number of single steps all pawns of $p$ still need to reach their last rank. A pawn move lowers $\pi$ (a step
advances a pawn; a capture advances the capturer and removes a pawn) and a heap move lowers $h$. We show, by induction on
$\pi(p)+h$, that for every node $p$ of $R(g)$ the player $X$ to move wins $p+\nim h$ iff $X$ wins $p+\nim h$ in $R(g)+\nim h$. A player wins iff
some option is lost by the opponent moving first. In $p+\nim h$ a T-move wins at once and an L-move loses at once. After a P-move to $x$
the opponent $Y$ wins $x+\nim h$ moving first by his terminal reply, so a P-move is never a winning option. A K-move leads to a node, and a
heap move to $p+\nim{h'}$, $h'<h$, with the board and flag unchanged; for both, by induction, the opponent loses moving first in the full
game iff he does in the reduced game. So $X$ has a winning option in one game iff in the other. Take $p=g$.
\end{proof}

\begin{lemma}[Wall pawns; \lab{P}]
Let $p$ have no passed pawn, with an e.p.\ flag (if any) protecting a pawn of the player not to move. Let $Z$ be the player to move and
$q$ a pawn of $Z$ on file $j$, where files $j-1,j,j+1$ are pattern files of one wall.
(a) After any move of $q$, $Z$ has no passed pawn.
(b) If moreover $p$ has no flag and files $j-1,\dots,j+3$ or files $j-3,\dots,j+1$ are pattern files of that wall, every move of $q$ is
L or P.
\end{lemma}
\begin{proof}
Files $j\pm1$ hold only an enemy pawn two ranks ahead of $q$. So the squares $q$ could capture on are empty. An e.p.\ capture by $q$
would need an enemy pawn beside $q$ on file $j\pm1$; there is none. File $j$ holds only $q$. Hence $q$ only advances: W3$\to$4, B5$\to$4 (type O),
W2$\to$3, W2$\to$4, B4$\to$3 (type E).
(a) The status of a pawn depends only on the enemy pawns on its three files and on the flag. No enemy pawn moves, and the old flag,
if any, protected an enemy pawn; so no other pawn of $Z$ becomes passed. After a single step $q$ still has an enemy pawn strictly
ahead on $j\pm1$. After W2$\to$4 the $B4$ pawns on $j\pm1$ stand beside $q$, so a flag is recorded and $q$ is not passed.
(b) Choose $s\in\{1,-1\}$ such that $j-s,j,j+s,j+2s,j+3s$ are pattern files. By (a) $X=Z$ has no passed pawn after the move. If $Y$
has one, the move is L. Otherwise $Y$ captures $q$ with his pawn on $j+s$: after a single step this pawn stands diagonally in front of $q$, and after
W2$\to$4 it stands beside $q$ and takes it en passant. Now files $j-1,j,j+1$ hold no pawn of $X$: file $j-s$ holds a pawn of $Y$ only, file
$j$ only the capturer, file $j+s$ nothing. So the capturer is passed (no flag after a capture). By (a) no pawn of $X$ was passed before the reply. A pawn of $X$ can change status only if a pawn of $Y$ on its three files left or
arrived, or a flag protecting it lapsed. The only flag (after W2$\to$4) protected $q$, which is gone; the arrival on $j$ only adds a
blocker; and file $j+s$ is one of the three files only of pawns on $j$, $j+s$, $j+2s$. Files $j$ and $j+s$ hold no $X$ pawn, and the
$X$ pawn on $j+2s$ is still blocked by the $Y$ pawn on $j+3s$. So $Y$ alone has a passed pawn, the reply wins for $Y$, and the move is P.
\end{proof}
The side $s$ matters: for the window $j-3,\dots,j+1$ the capture from $j+1$ can release the enemy pawn on $j+2$ and lose a race.

\paragraph{Insertion.}
Let $q$ be a position on $\cyl{n}$ and $0\le c\le n-2$, such that files $c,c+1$ are pattern files of one wall and the flagged pawn, if
any, is on neither of them. $\ins_c(q)$ on $\cyl{n+2}$ is $q$ with two new files placed after file $c$, holding copies of the contents
of files $c+1$ and $c$: old file $f$ goes to $f$ if $f\le c$ and to $f+2$ if $f\ge c+1$; a flag goes with its pawn, and the same side
holds the right. The new files continue the alternation, so they are pattern files of the lengthened wall. Two facts follow.
(V) The neighbours of every old file are copies of its old neighbours (old file $c$ now sees a copy of file $c+1$, old file $c+1$ a copy
of file $c$; the seam joins files $n+1$ and $0$), and the neighbours of the flagged pawn are the images of its old neighbours. So every old pawn has the same three files and
the same flag status: the same moves and the same status. A new pawn has an enemy pawn two ranks ahead on both neighbouring files and is
not passed. Hence the passed pawns of $\ins_c(q)$ are the images of those of $q$, with the same colour, rank, own-file content and
flag status; for the same last mover $q$ and $\ins_c(q)$ are both non-terminal, or both terminal with the same winner.
(C) If a pawn move $\nu$ in $q$ changes no square on files $c,c+1$, the image pawn has the corresponding move $\hat\nu$ in $\ins_c(q)$, and
$\ins_c(q\nu)=\ins_c(q)\hat\nu$. (Files $c,c+1$ of $q\nu$ are unchanged, a new flag lies on the file of $\nu$, and by (V) a double step
records a flag in both positions or in neither.)

\begin{lemma}[Insertion; \lab{P}]
Let $p$ on $\cyl{n}$ have no passed pawn and no e.p.\ flag, let $4\le c\le n-6$, and let files $c-4,\dots,c+5$ be pattern files of one
wall. Let $\hat p=\ins_c(p)$. Then $\hat p$ has no passed pawn and no flag, and for each side: the children of its K-moves in $\hat p$ are
the $\ins_c$-images of the children of its K-moves in $p$, and it has a T-move in $\hat p$ iff it has one in $p$.
\end{lemma}
\begin{proof}
By (V), $\hat p$ has no passed pawn; it has no flag. Files $c-4,\dots,c+7$ of $\hat p$ are pattern files. Let $X$ move.
(i) \emph{Central pawns.} A pawn of $X$ on a file $j\in\{c-3,\dots,c+4\}$ of $p$ has the window $j-1,\dots,j+3$ (if $j\le c+2$) or
$j-3,\dots,j+1$ (if $j\ge c-1$) inside $c-4,\dots,c+5$. By Lemma~2(b) all its moves are L or P. The same holds in $\hat p$ for the pawns on
files $c-3,\dots,c+6$ (the new pawns included), with windows inside $c-4,\dots,c+7$.
(ii) \emph{Far moves.} Every other pawn of $X$ is on a file $f\le c-4$ or $f\ge c+5$ of $p$, and its image is on $f$ or $f+2$ of $\hat p$;
every pawn of $\hat p$ not covered by (i) is such an image. A move $\mu$ of such a pawn changes only its own file and a neighbouring
file (the target of a capture, which is also the file of an e.p.\ victim), so only files $\le c-3$ or $\ge c+4$ (files modulo $n$), and a new
flag lies on its own file. Files $c-2,\dots,c+3$ of the child $p\mu$ are unchanged, so $\ins_c(p\mu)$ is defined and, by (C), equals $\hat p\hat\mu$.
By (V), $\mu$ and $\hat\mu$ are both T, both L, or both non-terminal.
(iii) \emph{Poison.} Let $p\mu$ be non-terminal. Its files $c-2,\dots,c+3$ are pattern files, and files $c-2,\dots,c+5$ of $\hat p\hat\mu$ are.
Call a pawn of $Y$ \emph{near} if it is on files $c-1,\dots,c+2$ of $p\mu$, resp.\ $c-1,\dots,c+4$ of $\hat p\hat\mu$, and \emph{far} otherwise.
A near pawn has three pattern files; the position has no passed pawn and its flag, if any, protects the pawn of $X$ that just moved. By
Lemma~2(a), after a near reply $Y$ has no passed pawn, so the reply does not win for $Y$; this holds on both boards.
The far pawns of $p\mu$ and of $\hat p\hat\mu$ correspond, with the same moves (V). A far reply $\nu$ starts on a file $\le c-2$ or $\ge c+3$
and changes only files $\le c-1$ or $\ge c+2$; an e.p.\ reply takes the pawn moved by $\mu$, which stands on a file $\le c-4$ or
$\ge c+5$, from a neighbouring file. So $\nu$ does not touch files $c,c+1$, and $\ins_c(p\mu\nu)=\hat p\hat\mu\hat\nu$ by (C). By (V),
$\nu$ wins for $Y$ iff $\hat\nu$ does. Hence $\mu$ is P iff $\hat\mu$ is.
By (i)--(iii), the T-moves of $X$ in $\hat p$ are the images of those in $p$, and the K-moves of $X$ in $\hat p$ are the images $\hat\mu$ of
the K-moves $\mu$ in $p$, with children $\ins_c(p\mu)$.
\end{proof}

\paragraph{Finite checks.} Let $n_0=17$ for odd $n$ and $n_0=16$ for even $n$. Let $Z_n=\{n-1,0,1,2,3,4\}$ (odd $n$) and
$Z_n=\{n-1,0,1,2,3\}$ (even $n$).
\begin{itemize}
\item[\textbf{C1}] \lab{C2} The reduced game $R(G_{n_0})$ has 144 nodes ($O_{17}$) resp.\ 104 nodes ($E_{16}$), with 209 resp.\ 177
K-edges; no node has an e.p.\ flag; every node agrees with $\mathrm{start}_{n_0}$ outside $Z_{n_0}$; and $R(G_{n_0})+\nim2$ is a
P-position. The graph (nodes and, for each side, the K-children and the existence of a T-move) was built by \texttt{rg.c} (C) and,
identically, by a second builder that uses the move generator, passed-pawn test and race rule of the library \texttt{cylcgt}; the
Python engine \texttt{eng.py} gives the same node and move counts, no flag and the same states on $Z_{n_0}$. The P-outcome at heap 2 was
obtained by solving this graph (\texttt{rg.c}, \texttt{eng.py}), by the canonical value $\nim2$ of the \texttt{cylcgt}-built reduced game
in \texttt{cylcgt}'s CGT arithmetic, and, without reduction, by C2 at $n_0$.
\item[\textbf{C2}] \lab{C2} For every odd $n$ with $5\le n\le61$ and every even $n$ with $6\le n\le60$, two independently written
unreduced whole-board searches, \texttt{rg.c} (C) and \texttt{eng.py} (Python), give the outcomes N, N, P, N for $G_n+\nim h$,
$h=0,1,2,3$, for each side moving first. Neither search uses the reduced game or Lemmas 1--3. In addition, the canonical value of the whole
board $G_n$, computed with \texttt{cylcgt}, is $\nim2$ for $5\le n\le10$. The proof uses only the widths $n<n_0$.
\end{itemize}

\begin{proof}[Proof of the Theorem]
Let $c=9$ (odd $n$) or $c=8$ (even $n$). In $\mathrm{start}_n$ every pawn has an enemy pawn strictly ahead on its own or a neighbouring
file, so it has no passed pawn; it has no flag; and for $n\ge n_0$, $\ins_c(\mathrm{start}_n)=\mathrm{start}_{n+2}$ (the copies of files
$c+1,c$ add one period to the wall). For $n\ge n_0$ of the right parity let $I_n$ say: no node of $R(G_n)$ has a flag, and every node agrees with
$\mathrm{start}_n$ outside $Z_n$. $I_{n_0}$ is C1.
Assume $I_n$. Every node $p$ is non-terminal (it is the root or the child of a K-move) and has no flag. The files $c-4,\dots,c+5$
($5,\dots,14$, resp.\ $4,\dots,13$) are wall files outside $Z_n$, because $n-2\ge14$ for odd $n\ge17$ and $n-2\ge13$ for even
$n\ge16$; so they hold their start content and are pattern files. Also $4\le c\le n-6$. So Lemma~3 applies to every node.
$\ins_c$ is injective (delete the two new files). By Lemma~3, the image of a K-edge is a K-edge, so, by induction on the distance from
$\mathrm{start}_n$, the images of the nodes of $R(G_n)$ are nodes of $R(G_{n+2})$; and every K-child of an image $\ins_c(p)$ is the image of a K-child of $p$, so, by induction on the distance
from $\mathrm{start}_{n+2}=\ins_c(\mathrm{start}_n)$, every node of $R(G_{n+2})$ is an image. Thus $\ins_c$ is an isomorphism of the
reduced games: nodes, side-labelled K-edges and, for each side, the existence of a T-move. Images have no flag. The new files are copies of files outside $Z_n$, $Z_n\setminus\{n-1\}$ lies left of the cut, and $n-1$ goes to
$n+1$; so every image agrees with $\mathrm{start}_{n+2}$ outside $Z_{n+2}$, and $I_{n+2}$ holds.
Hence $I_n$ holds and $R(G_n)\cong R(G_{n_0})$ for every $n\ge n_0$ of the right parity. The outcome of $R(g)+\nim h$ depends only on
this structure and $h$. So $R(G_n)+\nim2$ is a P-position by C1, and
$G_n+\nim2$ is a P-position by Lemma~1. For $n<n_0$ this is C2 (odd $5,\dots,15$; even $6,\dots,14$). No width is skipped.
Other heaps: in $G_n+\nim2$ every option loses, in particular the heap moves to $G_n+\nim0$ and $G_n+\nim1$ with the other side to
move; as this holds for either side moving first, $G_n+\nim0$ and $G_n+\nim1$ are N-positions. For $h\ge3$ the heap move to $G_n+\nim2$
wins, so $G_n+\nim h$ is an N-position. These are the outcomes of $\nim2+\nim h$.
(\lab{P}, given the computed facts C1 and C2.)
\end{proof}

\paragraph{What is and is not proved.}
The theorem gives the outcome of $G_n+\nim h$ for every heap; the proof uses no bound on $h$. Moreover, replace the child of every move
to a terminal position by the integer $+256$ (won by White) or $-256$ (won by Black). This gives a short partizan game $\hat G_n$, and
$G_n+\nim h$ has the outcome of the disjunctive sum $\hat G_n+\nim h$ (a heap move keeps the board and its flag, and $\pm256+\nim{h'}$ is won
by the favoured side whoever moves). As $\hat G_n+\nim2$ is a P-position, $\hat G_n=\nim2$ in normal-play CGT. This does not cover sums of
$G_n$ with other pawn positions when the first passed pawn on any board ends the game, nor the component-sum model of \texttt{cylcgt},
in which a ring that splits into components is evaluated as a sum; Lemma~1 is not available there. Only the race rule ``count'' is used.
\end{wallproof}

\bigskip\hrule\bigskip
\section*{Details (not part of the proof)}

\paragraph{Reachable positions.} By C1 and the induction, for every $n\ge n_0$ the reduced game has exactly 144 (odd) resp.\ 104 (even)
nodes. A node is $\mathrm{start}_n$ with the contents of the files of $Z_n$ replaced by one of the zone states listed in the code
repository. $Z_n$ holds files $0,1,2$ and the wall pawns next to them (odd: $B5$ on $n-1$, $W3$ on 3, $B5$ on 4; even: $W2$ on $n-1$ and
on 3). In every node of the reduced game the wall files $5,\dots,n-2$ (odd) resp.\ $4,\dots,n-2$ (even) keep their start content. This
says nothing about play after an L- or P-move, or after a T-move.

\paragraph{The five poisoned wall moves (Lemma~2(b)).}
The table lists the wall advances and the poisoning reply of $Y$ when the advance does not already lose (if $Y$ has a passed pawn after
it, the move is L). It shows the window $j-1,\dots,j+3$ ($s=1$); for the window $j-3,\dots,j+1$ replace $j+k$ by $j-k$.
\begin{center}\small
\begin{tabular}{llll}
type & move of $q$ on file $j$ & reply of $Y$ & why $X$ has no passed pawn afterwards\\\hline
O & W3$\to$4 & B5 on $j{+}1$ takes on $(j,4)$ & W3 on $j{+}2$ blocked by B5 on $j{+}3$\\
O & B5$\to$4 & W3 on $j{+}1$ takes on $(j,4)$ & B5 on $j{+}2$ blocked by W3 on $j{+}3$\\
E & W2$\to$3 & B4 on $j{+}1$ takes on $(j,3)$ & W2 on $j{+}2$ blocked by B4 on $j{+}3$\\
E & W2$\to$4 & B4 on $j{+}1$ takes e.p.\ on $(j,3)$ & same\\
E & B4$\to$3 & W2 on $j{+}1$ takes on $(j,3)$ & B4 on $j{+}2$ blocked by W2 on $j{+}3$\\
\end{tabular}
\end{center}
In each case the capturer is passed because files $j-1,j,j+1$ then hold no pawn of $X$. Example for the side: on eleven files with
$W3,B5,W3,B5,W3,B5$ on files $3,\dots,8$ and nothing else, B5$\to$4 on file 6 has only the window $3,\dots,7$. The capture from file 5
wins for White; the capture from file 7 releases the $B5$ on file 8, both distances are 4, and Black wins the tie.

\paragraph{Remarks.}
(1) Only the whole-board game is covered. In the component-sum model of \texttt{cylcgt}, Lemma~1 is not available: a decisive passer in one
component can be cancelled by a $-256$ option in another.
(2) The ten-file pattern interval (radius 5) in Lemma~3 is what the hand proof uses; it is a sufficient hypothesis. In a random sample
for the odd family, the correctly phased insertion had counterexamples at radius 1 and 2 and none at radius 3, and a wrongly phased
insertion had counterexamples at radius 5; the correctly phased samples at radius 5 had none.
(3) The same induction applies to the single-step family with core $K=(W3B6,B6,W4B7,W4,W6B7)$: $F_n=(B4,W2)^{(n-5)/2}K$ for odd
$n\ge5$ and $F_n=(B4,W2)^{(n-6)/2}K\,W2$ for even $n\ge8$. Its reduced games have 662 nodes (even $n$, $n_0=18$, cut $c=4$, zone
$\{n-7,\dots,n-3\}$) and 3568 nodes (odd $n$, $n_0=19$, cut $c=6$, zone $\{0,1,n-6,\dots,n-2\}$); in both the ten-file interval lies in the
unchanged part of the wall. The analogues of C1 and C2 were computed for it (C2 by the same two searches for $n=5$ and $7\le n\le61$).
This transfer has not been written out in full.
\end{document}
